%|||||||||||||||||||||||||||||||*******************************|||||||||||||||||||||||||||||||||||||%
%|||||||||||||||||||||||||||||||*******************************|||||||||||||||||||||||||||||||||||||%
%|||||||||||||||||||||||||||||||*******************************|||||||||||||||||||||||||||||||||||||%
%|||||||||||||||||||||||||||||||********KUMPULAN-SCRIPT********|||||||||||||||||||||||||||||||||||||%
%|||||||||||||||||||||||||||||||*******************************|||||||||||||||||||||||||||||||||||||%
%|||||||||||||||||||||||||||||||*******************************|||||||||||||||||||||||||||||||||||||%
%|||||||||||||||||||||||||||||||*******************************|||||||||||||||||||||||||||||||||||||%

% Membuat Frame Menjadi 2
% ---------------------------------------------------------------------
\begin{frame}
  \frametitle{Frame Dibagi Menjadi Dua}

  \begin{columns} %Frame Kiri
    \begin{column}{0.5\textwidth}
      \justifying % Perataan kanan kiri
      \tiny % Ukuran teks menjadi kecil
      Konten kolom pertama dengan perataan kiri.
    \end{column}
    
    \begin{column}{0.5\textwidth} % Frame Kanan
      \justifying % Perataan kanan kiri
      \tiny % Ukuran teks menjadi kecil
      Konten kolom kedua dengan perataan kanan dan ukuran teks besar.
    \end{column}
  \end{columns}
\end{frame}
% ---------------------------------------------------------------------

% ITEMIZE BERBENTUK PANAH
% ---------------------------------------------------------------------
\begin{itemize}
    \item[$\Rightarrow$]
\end{itemize}
% ---------------------------------------------------------------------

% 1. Upright Roman font : Its alphabets are typeset using the command \mathrm{}.

% 2. Normal math italic font : for typesetting normal mathematics letters. The command is \mathnormal{}. In summary, {$\mathnormal{a}$} is the same as $a$.

% 3.  Calligraphic font : for typesetting capital letters with a special Calligraphic font. The command here is \mathcal{}.

% 4. Upright Roman boldface : for typesetting upright roman boldface letters. The command here is \mathbf{}.

% 5. Upright sans serif : for typesetting upright sans serif letters The command here is \mathsf{}.

% 6. Italic letters : for typesetting text in italic letters. The command here is \mathit{}.

% 7. Typewriter type font : for typesetting upright typewriter type letters. The command here is \mathtt{}.\

% Tulisan Kecill
\fontsize{4}{6}